\documentclass[11pt]{article}
\usepackage[margin=1in]{geometry}
\usepackage{amsmath,amssymb,amsthm}
\usepackage{graphicx}
\usepackage{microtype}
\usepackage{xcolor}
\usepackage[colorlinks=true,linkcolor=blue!55!black,citecolor=blue!55!black,urlcolor=blue!55!black]{hyperref}

\newtheorem{theorem}{Theorem}
\newtheorem{lemma}{Lemma}
\newcommand{\K}{\mathbb K}
\newcommand{\N}{\mathbb N}
\newcommand{\Lip}{\mathrm{Lip}}
\newcommand{\rank}{\operatorname{rank}}

\title{Coordinatewise products on $\ell_1$ are Lipschitz $p$-summing}
\author{Candidate full solution packet for arXiv:1805.02115}
\date{22 July 2026}

\begin{document}
\maketitle

\begin{center}
\fbox{\parbox{0.91\textwidth}{\textbf{Status.} Candidate full solution,
likely valid, pending human review.  The theorem below answers the source
question positively for every $1\le p<\infty$ and proves an arity-independent
strengthening.}}
\end{center}

\section{Source question}

Angulo-L\'opez and Fern\'andez-Unzueta introduce Lipschitz $p$-summing
multilinear operators in \cite{AnguloFernandez2020}.  On page 6, immediately
after Proposition 3.5, they ask whether the coordinatewise product
\[
 T:\ell_1\times\ell_1\longrightarrow\ell_1,
 \qquad T(x,y)=(x_jy_j)_{j\ge1},
\]
is Lipschitz $p$-summing.  Here is the exact source passage.

\begin{figure}[ht]
\centering
\includegraphics[width=\textwidth]{figures/open_problem_crop.png}
\caption{Source page 6: the sentence posing the coordinatewise-product
question, with the preceding difference-tensor context.}
\end{figure}

For completeness, an $m$-linear map $S:X_1\times\cdots\times X_m\to Y$ is
Lipschitz $p$-summing if a constant $C$ exists such that, for every finite
family $u_i,v_i\in X_1\times\cdots\times X_m$,
\begin{equation}\label{eq:def}
 \left(\sum_i\|S(u_i)-S(v_i)\|^p\right)^{1/p}
 \le C\sup_{\varphi\in B_{\mathcal L(X_1,\ldots,X_m)}}
 \left(\sum_i|\varphi(u_i)-\varphi(v_i)|^p\right)^{1/p}.
\end{equation}
The least such $C$ is denoted $\pi_p^{\Lip}(S)$.

\section{Proof intuition}

The definition only tests differences of two decomposable tensors.  After
flattening one tensor coordinate against all the others, such a difference is
a matrix of rank at most two.  Its coordinate diagonal has $\ell_1$ norm at
most the matrix trace norm, and rank two bounds that trace norm by $\sqrt2$
times the Hilbert--Schmidt norm.

The missing bridge to the right-hand side of \eqref{eq:def} is supplied by
random signs.  A bounded sign array is a norm-one multilinear form on a product
of $\ell_1$ spaces.  The lower Khintchine inequality recovers the
Hilbert--Schmidt norm of each difference array from the $L_p$ norm of its
random-sign evaluation.  Crucially, one first sums over the finite family and
only then averages; the expectation is bounded by one common sign-array
supremum, precisely the supremum allowed in \eqref{eq:def}.

\section{The result}

Let $\kappa_{p,\K}>0$ be any lower scalar Khintchine constant over
$\K\in\{\mathbb R,\mathbb C\}$, so that for every square-summable scalar
family $(a_\alpha)$ and independent Rademacher signs $(\varepsilon_\alpha)$,
\begin{equation}\label{eq:khin}
 \left(\mathbb E\left|\sum_\alpha
 \varepsilon_\alpha a_\alpha\right|^p\right)^{1/p}
 \ge \kappa_{p,\K}\left(\sum_\alpha|a_\alpha|^2\right)^{1/2}.
\end{equation}
Only positivity of this standard constant is used; see
\cite{Khintchine1923,Kahane1985}, including the complex-scalar formulation.

\begin{theorem}\label{thm:main}
Let $m\ge2$ and $1\le p<\infty$.  Over either scalar field, the coordinatewise
product
\[
 M_m:\underbrace{\ell_1\times\cdots\times\ell_1}_{m\ \mathrm{times}}
 \longrightarrow\ell_1,
 \qquad
 M_m(x^1,\ldots,x^m)=\bigl(x^1_j\cdots x^m_j\bigr)_{j\ge1},
\]
is Lipschitz $p$-summing, and
\[
 1\le \pi_p^{\Lip}(M_m)\le \frac{\sqrt2}{\kappa_{p,\K}}.
\]
In particular, the question on page 6 of \cite{AnguloFernandez2020} has an
affirmative answer for every $1\le p<\infty$.
\end{theorem}

We isolate the rank-two observation.

\begin{lemma}\label{lem:diag}
Suppose
\[
 z=x^1\otimes\cdots\otimes x^m-y^1\otimes\cdots\otimes y^m
 \in \ell_1(\N^m),
\]
and set $\Delta z=(z_{j,\ldots,j})_{j\ge1}$.  Then
\begin{equation}\label{eq:diag}
 \|\Delta z\|_1\le \sqrt2\left(\sum_{\mathbf j\in\N^m}
 |z_{\mathbf j}|^2\right)^{1/2}.
\end{equation}
\end{lemma}

\begin{proof}
Regard $z$ as the matrix of an operator
\[
 Z:\ell_2(\N^{m-1})\longrightarrow\ell_2(\N)
\]
by placing the first index in the row and the remaining indices in the column.
Each decomposable tensor gives a rank-one operator under this flattening, so
$\rank Z\le2$.  Moreover, $Z$ is Hilbert--Schmidt and
\[
 \|Z\|_{S_2}=\left(\sum_{\mathbf j}|z_{\mathbf j}|^2\right)^{1/2}.
\]
Choose scalars $\theta_j$ of modulus at most one so that
$\theta_jz_{j,\ldots,j}=|z_{j,\ldots,j}|$, and let
$Q:\ell_2(\N)\to\ell_2(\N^{m-1})$ be the diagonal partial isometry
$Qe_j=\theta_j e_{(j,\ldots,j)}$.  Then $\|Q\|\le1$, and trace duality gives
\[
 \|\Delta z\|_1
 =|\operatorname{tr}(ZQ)|
 \le \|Z\|_{S_1}\|Q\|
 \le \|Z\|_{S_1}.
\]
If $s_1,s_2$ are the only possibly nonzero singular values of $Z$, then
\[
 \|Z\|_{S_1}=s_1+s_2
 \le\sqrt2(s_1^2+s_2^2)^{1/2}
 =\sqrt2\|Z\|_{S_2},
\]
which proves \eqref{eq:diag}.  The argument applies directly to the infinite
array because $z\in\ell_1(\N^m)\subset\ell_2(\N^m)$ and $Z$ has finite rank.
\end{proof}

\begin{proof}[Proof of Theorem \ref{thm:main}]
For a sign array $\varepsilon=(\varepsilon_{\mathbf j})_{\mathbf j\in\N^m}$,
define
\[
 \varphi_\varepsilon(x^1,\ldots,x^m)
 :=\sum_{\mathbf j=(j_1,\ldots,j_m)\in\N^m}
 \varepsilon_{\mathbf j}x^1_{j_1}\cdots x^m_{j_m}.
\]
The series is absolutely convergent, and
\[
 |\varphi_\varepsilon(x^1,\ldots,x^m)|
 \le \prod_{r=1}^m\|x^r\|_1.
\]
Because every coefficient has modulus one, $\|\varphi_\varepsilon\|=1$.

Now fix finitely many pairs
$u_i=(x_i^1,\ldots,x_i^m)$ and
$v_i=(y_i^1,\ldots,y_i^m)$, and put
\[
 z_i=x_i^1\otimes\cdots\otimes x_i^m
     -y_i^1\otimes\cdots\otimes y_i^m.
\]
The output difference is $M_m(u_i)-M_m(v_i)=\Delta z_i$.  Lemma
\ref{lem:diag} and \eqref{eq:khin} yield
\[
 \|M_m(u_i)-M_m(v_i)\|_1^p
 \le \left(\frac{\sqrt2}{\kappa_{p,\K}}\right)^p
 \mathbb E\,|\langle\varepsilon,z_i\rangle|^p.
\]
Summing over $i$ before taking the supremum gives
\begin{align*}
 \sum_i\|M_m(u_i)-M_m(v_i)\|_1^p
 &\le \left(\frac{\sqrt2}{\kappa_{p,\K}}\right)^p
 \mathbb E\sum_i
 |\varphi_\varepsilon(u_i)-\varphi_\varepsilon(v_i)|^p\\
 &\le \left(\frac{\sqrt2}{\kappa_{p,\K}}\right)^p
 \sup_{\varphi\in B_{\mathcal L(\ell_1,\ldots,\ell_1)}}
 \sum_i|\varphi(u_i)-\varphi(v_i)|^p.
\end{align*}
After taking $p$th roots, this is \eqref{eq:def} with
$C=\sqrt2/\kappa_{p,\K}$.  Finally,
$\pi_p^{\Lip}(M_m)\ge\|M_m\|=1$, giving the stated two-sided bound.
\end{proof}

\section{Verification and limitations}

The included script checks 960 real/complex decomposable differences of
arities two and three, verifies rank at most two and
$\|\Delta z\|_1\le\sqrt2\|z\|_2$, and performs exact finite sign enumeration
for small families.  The code is only a sanity check; the proof above is
dimension-free and analytic.

This packet does not claim the exact value of $\pi_p^{\Lip}(M_m)$.  It also
does not claim the separate Section 7 question of \cite{AnguloFernandez2020}:
Fern\'andez-Unzueta explicitly answered that question in
\cite{Fernandez2022}.  A bounded novelty search used the run indexes, exact
arXiv phrase/title queries, and the local full text of later arXiv papers
citing the source.  No prior answer to the coordinatewise-product question or
the rank-two/random-sign argument was found.  Novelty confidence is moderate
pending specialist review.

\section{Human-review recommendation}

The candidate should receive expert proof review.  The highest-value checks
are: (i) the trace pairing with the diagonal partial isometry in Lemma
\ref{lem:diag}; (ii) the scalar Khintchine inequality over the complex field;
and (iii) the order ``sum, then average, then supremum'' in the final display.

\begin{thebibliography}{9}

\bibitem{AnguloFernandez2020}
J.~C. Angulo-L\'opez and M. Fern\'andez-Unzueta,
\emph{Lipschitz $p$-summing multilinear operators},
J. Funct. Anal. 279 (2020), no.~4, 108572;
\href{https://arxiv.org/abs/1805.02115}{arXiv:1805.02115}.

\bibitem{Fernandez2022}
M. Fern\'andez-Unzueta,
\emph{Lipschitz $p$-summing multilinear operators correspond to Lipschitz
$p$-summing operators},
\href{https://arxiv.org/abs/2204.01775}{arXiv:2204.01775} (2022).

\bibitem{Khintchine1923}
A. Khintchine,
\emph{\"Uber dyadische Br\"uche},
Math. Z. 18 (1923), 109--116.

\bibitem{Kahane1985}
J.-P. Kahane,
\emph{Some Random Series of Functions}, second ed.,
Cambridge Studies in Advanced Mathematics 5, Cambridge University Press, 1985.

\end{thebibliography}

\end{document}
